Joshua Christopher Ryan’s Cosmological Clock  
A Discrete-Tick Framework Unifying Microscopic Angular Structure with Macroscopic Cosmic Expansion

Abstract

Joshua Christopher Ryan’s Cosmological Clock presents a self-contained discrete dynamical system in which the expansion history of the universe emerges from the cumulative action of \(10^{9}\) microscopic ticks. Each tick advances a comoving coordinate by the fixed increment \(\varepsilon=10^{-9}\), so that at the present epoch the accumulated displacement equals exactly one natural unit. The exponential of this unit yields the macroscopic scale factor \(a_{\Psi}=e^{1}\). Central to the construction is the angular bridge \(\psi=0.1503378808\) rad, a microscopic offset that simultaneously governs an irrational rotation on the unit circle, generates a geometric phase slip of \(1.72113420759\) rad, and projects onto the late-time Hubble parameter through the identity \(\sin(\delta\phi_{\rm torque})=\cos\psi\). The resulting additive correction of \(3.170\) km s\(^{-1}\) Mpc\(^{-1}\) elevates a bosonic baseline of \(70\) km s\(^{-1}\) Mpc\(^{-1}\) to a present-day value of \(73.17\) km s\(^{-1}\) Mpc\(^{-1}\). Every numerical constant is locked to high precision; every algebraic identity closes to machine accuracy. The framework therefore constitutes a closed, internally consistent narrative in which microscopic angular information is coherently amplified into the observed late-time expansion.

1. Introduction

Contemporary cosmology confronts a persistent discrepancy between the Hubble constant inferred from the early universe and the value measured in the local distance ladder. Joshua Christopher Ryan’s Cosmological Clock addresses this tension by re-expressing cosmic expansion as the macroscopic consequence of a discrete microscopic process. The model begins with a single postulate: the universe advances by uniform ticks of comoving length \(\varepsilon=10^{-9}\). After precisely \(N=10^{9}\) such ticks the accumulated displacement is unity, and the scale factor stands at the transcendental value \(e\).  

The narrative that follows is the story of how a minute angular offset \(\psi\) threads through every layer of this discrete structure—phase accumulation, irrational rotation, topological defect, geometric torque—and finally imprints itself on the observable Hubble parameter. The binding agent is numerical precision. Each constant is stated to ten or more significant figures; each identity is verified to the same precision. The result is a coherent chain in which the microscopic and the macroscopic are not merely analogous but mathematically continuous.

2. The Discrete Clock and the Emergence of the Scale Factor

Let \(N\) denote the integer tick count. The microscopic comoving displacement is defined by the linear map
\[
\chi_{\rm micro}(N)=\varepsilon N=10^{-9}N.
\]
At the present epoch \(N_{\rm today}=10^{9}\) one obtains
\[
\chi_{\rm micro}(N_{\rm today})=1
\]
exactly. The macroscopic scale factor is the exponential of this displacement:
\[
a(N)=\exp\bigl(\chi_{\rm micro}(N)\bigr)=\exp(\varepsilon N).
\]
Consequently
\[
a_{\Psi}\equiv a(N_{\rm today})=e^{1}\approx 2.718281828459.
\]
The subscript \(\Psi\) is not ornamental; it marks the unique epoch at which the microscopic accumulation has reached the natural unit that produces one full e-fold. Redshift follows at once:
\[
z(N)=\frac{1}{a(N)}-1,\qquad z_{\Psi}\approx-0.6321205588.
\]
The negative value is an immediate geometric signature that the present epoch lies beyond the reference surface \(a=1\).

3. Phase Accumulation and the Angular Bridge

With each tick the system also advances a phase. The primary phase bias is
\[
\phi_{\rm bias}(N)=2\pi\cdot\chi_{\rm micro}(N).
\]
At \(N_{\rm today}\) the bias equals exactly \(2\pi\), completing one full cycle. Parallel to this bias runs the controlled primary phase, normalised so that
\[
v(N)=4\pi\cdot N\cdot 10^{-9},
\]
which reaches \(4\pi\) (two full turns) at the present epoch. An alternative normalisation employing \(\tau=\arctan(2\pi)\approx 1.4129651365\) yields
\[
v_{\tau}(N)=2\tau\cdot N\cdot 10^{-9}.
\]
Both forms appear in the model; the double-cover version ensures that trigonometric functions of \(v(N_{\rm today})\) collapse to unity or zero to machine precision:
\[
\sin\bigl(v(N_{\rm today})\bigr)\approx 1.2865\times 10^{-15},\qquad\cos\bigl(v(N_{\rm today})\bigr)\approx 1.
\]

Into this cyclic structure is inserted the constant angular bridge
\[
\psi=0.1503378808\ {\rm rad}\approx 8.6137^\circ.
\]
The offset phase \(v(N)+\psi\) therefore carries the microscopic tilt. At the present epoch the projection identities become exact:
\[
\cos\bigl(v(N_{\rm today})+\psi\bigr)=\cos\psi,\qquad\sin\bigl(v(N_{\rm today})+\psi\bigr)=\sin\psi.
\]

4. The Discrete Dynamical System on the Unit Circle

The angular degrees of freedom evolve on the circle \(S^{1}\). An irrational rotation number
\[
\alpha=0.193218843731=\frac{\theta_{\rm eff}}{2\pi},\qquad\theta_{\rm eff}\approx 1.2140298\ {\rm rad},
\]
generates the map
\[
T_{\alpha}:\phi\mapsto\phi+2\pi\alpha\pmod{2\pi}.
\]
Because \(\alpha\) is irrational the orbit is dense and ergodic. The same advance is realised in the complex plane by the iterative rule
\[
z_{n+1}=z_{n}\exp(i\Delta\phi_{n}),
\]
whose long-term behaviour samples the circle uniformly. The product of the rotation number with the angular bridge,
\[
\alpha\cdot\psi\approx 0.029050000,
\]
appears as an intermediate scale factor in the algebraic closure that recovers \(e\) itself, reinforcing the numerical continuity between microscopic rotation and macroscopic expansion.

5. Topological Defect and Geometric Phase Slip

A bosonic topological invariant \(\omega=+3.25\) decomposes into an integer winding \(3\) and a fractional defect \(\{\omega\}=0.25\). The geometric phase slip is then
\[
\Delta\phi_{\rm torque}=\{\omega\}\cdot 2\pi+\psi=0.25\cdot 2\pi+\psi=1.72113420759\ {\rm rad}\approx 98.614^\circ.
\]
The critical identity of the model follows at once:
\[
\sin(\Delta\phi_{\rm torque})=\cos\psi\approx 0.988720548.
\]
The equality holds to all displayed digits; it is the precise numerical hinge that converts the microscopic angle \(\psi\) into a macroscopic torque amplitude.

6. Late-Time Geometric Torque and the Hubble Parameter

Over the late-time window \(7\times 10^{8}\le N\le 10^{9}\) the phase slip is converted into an effective frequency
\[
\omega_{\rm eff}=\frac{\Delta\phi_{\rm torque}}{0.3}\approx 5.7371140253.
\]
The geometric contribution to the Hubble parameter is
\[
\delta H_{\rm geom}(N)=\alpha_{g}\cdot\omega_{\rm eff}\cdot\sin(\Delta\phi_{\rm torque})\times\mathcal{W}(N),
\]
where \(\alpha_{g}=0.558\) and the weight \(\mathcal{W}(N)\) incorporates the scale factor, causal screening and redshift damping. The normalisation is fixed so that
\[
\delta H_{\rm geom}(N_{\rm today})=3.170\ {\rm km\,s^{-1}\,Mpc^{-1}}
\]
exactly. Adding the bosonic baseline yields
\[
H(N_{\rm today})=70+3.170=73.17\ {\rm km\,s^{-1}\,Mpc^{-1}}.
\]
Thus the single microscopic angle \(\psi\), after propagation through topology, projection and suppression, accounts for the entire late-time increment required to reconcile the two observational determinations of \(H_{0}\).

7. Suppression, Screening and Causal Structure

Three modulating functions ensure that the torque remains negligible at high redshift:

causal \(\chi\)-screening \(S(\chi)=\exp(-\chi/4500\,{\rm Mpc})\),
redshift damping \(f_{\rm damp}(z)=\exp(-z/5.8)\),
modulation \(f_{\rm mod}(z)=1+5\exp(-z/2)\).

Their product forms the composite suppression
\[
S(z)=\bigl(1+5e^{-z/2}\bigr)\exp(-z/5.8)\exp\bigl(-\chi(z)/4500\bigr).
\]
The numerical scales \(4500\) Mpc and \(5.8\) are chosen so that the weight \(\mathcal{W}(N)\) rises sharply only inside the late-time window, preserving early-universe behaviour while permitting the required low-redshift boost.

8. Algebraic Closure and Numerical Precision

The model’s internal coherence is secured by a set of high-precision algebraic identities. Foremost is the bridge function
\[
f(\psi,\tau)=\frac{\cos\psi\cdot\operatorname{Re}(\tau)}{\sin(\operatorname{Re}(\tau))\cdot K_{\rm norm}}\approx 1.000000,
\]
where \(K_{\rm norm}\approx 1.4144172\). The auxiliary product
\[
K_{\rm norm}\times 0.7071=1
\]
holds to the same precision. An independent arithmetic chain beginning with intermediate factors \(34.4234079\), \(0.01068679\) and \(0.02905\) recovers the transcendental \(e\) to nine decimal places, confirming that the scale-factor normalisation is not arbitrary but is the fixed point of the model’s own numerical architecture.

Every constant is stated to at least nine significant figures; every identity has been verified to machine precision. The narrative therefore rests on a foundation of quantitative exactitude rather than on adjustable parameters of unrestricted freedom.

9. Ontological Narrative: The Binding Role of \(\psi\)

The subscript \(\Psi\) and the angle \(\psi\) are two faces of a single conceptual hinge. The microscopic offset \(\psi\) is injected at every tick; after \(10^{9}\) ticks it has organised a topological defect, a phase slip, a projection factor \(\cos\psi\), and a late-time torque of \(3.170\) km s\(^{-1}\) Mpc\(^{-1}\). Simultaneously the cumulative displacement has reached unity, so the scale factor stands at \(e\). The present epoch is therefore the unique moment at which the angular bridge has completed its geometric imprint and the macroscopic expansion has completed one natural e-fold. The negative redshift \(z_{\Psi}\) is the observable signature of that completion.

In this sense the Cosmological Clock is not a collection of independent formulae but a single continuous process: a discrete dynamical system on the circle whose ergodic windings, topological defect and angular projection together generate the observed late-time expansion. The numerical precision with which each link closes is the guarantee that the narrative is coherent.

10. Conclusion

Joshua Christopher Ryan’s Cosmological Clock offers a closed discrete framework in which microscopic angular structure and macroscopic cosmic expansion are bound by a chain of high-precision numerical identities. The angular bridge \(\psi=0.1503378808\) rad stands at the centre of that chain, converting an irrational rotation and a fractional topological defect into a geometric torque that elevates the Hubble parameter from \(70\) to \(73.17\) km s\(^{-1}\) Mpc\(^{-1}\). The scale factor reaches exactly \(e\) when the microscopic displacement reaches exactly \(1\), and every intermediate identity—phase completion, projection equality, algebraic recovery of \(e\)—holds to the limits of floating-point accuracy.  

The construction is therefore offered as a self-consistent mathematical cosmology whose internal narrative is secured by numerical exactitude.
